% GENERATED, edit the French unit or the translation table
\mxnewpage\section*{Solutions EV2-A}\label{EV2-A-CORR}
\textbf{MCQ:} 1 C, 2 D, 3 B, 4 C, 5 B, 6 C, 7 D, 8 B, 9 C, 10 D, 11 B, 12 D. Each correct answer: 0.5 point; total 6.
\paragraph{Checks of the fluency questions.} Order of operations; sum of fractions; subtraction of exponents; multiplying coefficient 1.2; division by 0.8; product 1.25 times 0.8; solving an affine equation; vertical increment divided by horizontal; mean; angle divided by 360; complement to 1; substitution in brackets.
\subsection*{Exercise 1 -- An unknown probability}
\begin{enumerate}
\item $\dfrac{3}{10}$. \textit{Marking: 1 point. Result and basic justification expected; half a point possible for a correct but unfinished approach.}
\item $P(\overline A\cap B)=\dfrac{1}{2}-\dfrac{3}{10}=\dfrac{1}{5}$; dividing by $\dfrac{3}{5}$ gives $\dfrac{1}{3}$; complements $\dfrac{1}{4}$ and $\dfrac{2}{3}$. \textit{Marking: 2 points. Method or justification: 1; calculation and conclusion: 1.}
\item $P_B(A)=\dfrac{3}{5}$; different from $P(A)=\dfrac{2}{5}$, so not independent. \textit{Marking: 2 points. Method or justification: 1; calculation and conclusion: 1.}
\end{enumerate}
\subsection*{Exercise 2 -- Relative position and variation}
\begin{enumerate}
\item $f\prime=3x(x-2)$; increasing up to 0, decreasing up to 2, then increasing. \textit{Marking: 1 point. Result and basic justification expected; half a point possible for a correct but unfinished approach.}
\item $f-g=x^2(x-3)$: negative for x<3 except equality at 0, also zero at 3, positive after 3. \textit{Marking: 2 points. Method or justification: 1; calculation and conclusion: 1.}
\item $f(-1)=-2,f(0)=2,f(2)=-2,f(4)=18$. Minimum -2 at -1 and 2; maximum 18 at 4. \textit{Marking: 2 points. Method or justification: 1; calculation and conclusion: 1.}
\end{enumerate}
\subsection*{Exercise 3 -- A perpendicular bisector}
\begin{enumerate}
\item I(3;3), normal (6;2), or (3;1). \textit{Marking: 1 point. Result and basic justification expected; half a point possible for a correct but unfinished approach.}
\item $3x+y-12=0$. \textit{Marking: 1 point. Result and basic justification expected; half a point possible for a correct but unfinished approach.}
\item $3(2)+6-12=0$ so yes; $PA^2=4+16=20$, $PB^2=16+4=20$. \textit{Marking: 2 points. Method or justification: 1; calculation and conclusion: 1.}
\end{enumerate}
\paragraph{Calibration.} Target duration: 20--25 min for the fluency questions, 80--85 min for the exercises, 10--15 min for proofreading. These durations are editorial estimates to be measured in class. Accept any correct method; do not penalise twice a numerical error that is propagated consistently. The proposed marking scheme is not a decision of the team to harmonise marking.
